% Constraint-Enriched Decoupled Planning
% 中文初稿，英文SCI叙事逻辑，后续翻译
\documentclass[a4paper,11pt]{article}
\usepackage[UTF8]{ctex}
\usepackage{geometry,amsmath,amssymb,booktabs,graphicx,enumitem,algorithm2e}
\geometry{margin=2.5cm}
\setlist{nosep}

\title{约束增强的解耦规划框架：弥合多障碍物自动驾驶中的路径--速度信息断裂}
\author{廖嘉旺，孙成娇*\\湖北文理学院 计算机工程学院}
\date{}

\begin{document}
\maketitle

% ============================================================
\begin{abstract}
路径--速度解耦是量产自动驾驶规划器的主流架构，其将三维轨迹优化分解为两个顺序求解的二维子问题以降低计算复杂度。然而，这种串行分解在路径阶段丢弃了时间维度信息，导致动态障碍物场景下的结构性失效：路径阶段无法感知障碍物的时序行为，速度阶段则只能在固定路径上被动制动。本文提出约束增强解耦规划框架（Constraint-Enriched Decoupled Planning, CEDP），在保留解耦架构计算优势的前提下，通过建立路径阶段到速度阶段的显式约束传递通道恢复时空协调能力。该通道由三个协同组件构成：多因子威胁度过滤器确定需要传递约束的障碍物子集；最大间隙空间决策以$O(n\log n)$复杂度输出最优通行带；时空可通行走廊将到达时间约束注入速度规划的二次规划问题。理论分析给出了纯解耦规划器的不可行性充分条件，并证明当走廊非空时约束增强规划器的可行性保证。在Apollo 11.0生产级代码上的实现表明，框架仅新增约800行C++代码。15个仿真场景和实车验证结果表明，CEDP在多障碍物压力场景下将通过率从0\%提升至100\%，同时平均巡航速度提高29.3\%，规划延迟低于10\,ms，满足实时性要求。
\end{abstract}

\textbf{关键词：} 自动驾驶；运动规划；路径--速度解耦；约束传递；时空走廊；Apollo

% ============================================================
\section{引言}
% ============================================================

路径--速度解耦规划是当前量产自动驾驶系统中应用最广泛的轨迹规划架构。百度Apollo的EM Planner将轨迹优化分解为Frenet坐标系下的横向路径规划（SL域）和纵向速度规划（ST域），通过迭代求解两个低维二次规划问题来逼近联合最优解。京东配送车ROVER 5.0、Autoware的Lattice Planner等工业系统均采用类似的分解策略。这种架构的优势在于将$O(n^3)$级别的三维联合优化降解为两个$O(n)$级别的子问题，使规划周期稳定在10\,ms量级，满足实时性需求。

然而，路径--速度解耦在多障碍物动态场景下存在结构性信息断裂。路径阶段仅基于障碍物的空间投影做横向决策，不具备时间维度的感知能力；速度阶段接收固定路径后，面对动态障碍物只能选择制动或等待，无法通过调整横向位置来规避。当多个动态障碍物的时序行为与自车到达时间存在耦合关系时，这种信息断裂可导致规划器在联合空间本有可行解的情况下报告无解。

针对这一问题，现有研究主要沿两条路线展开。第一条路线是放弃解耦架构，转向时空联合规划。Yang等人提出IA Planner，通过瞬时分析将三维时空约束投影为二维空间约束后用MPC求解联合轨迹。胡杰等人基于改进Hybrid A*在三维$s$-$l$-$t$空间中搜索粗粒度轨迹，再用NMPC细化。这类方法以计算代价翻倍（76--117\,ms）为代价换取时空一致性。第二条路线是在解耦框架内做局部改进。Han等人在路径阶段的QP目标函数中注入障碍物预测概率作为soft cost，隐式传递预测信息。Zhang等人通过并行枚举候选路径--速度组合并用合作博弈选取最优配对，但其$O(mn)$枚举复杂度与架构改动量均较大。

本文提出第三条路线：在严格保留解耦架构的前提下，通过建立路径阶段到速度阶段的\textbf{显式约束传递通道}恢复时空协调能力。与soft cost隐式传递不同，本文的通道以硬约束形式将空间决策信息结构化地注入速度规划子问题，具有明确的可行性保证。与联合规划不同，通道仅在两阶段接口处增加信息，不改变各阶段的内部求解结构，计算代价几乎不增加。

本文的贡献如下：
\begin{enumerate}
  \item 提出约束增强解耦规划框架CEDP，通过三维度协同的显式通道（威胁度过滤$\to$空间决策$\to$时间约束注入）弥合路径--速度信息断裂；
  \item 给出纯解耦规划器的不可行性充分条件及约束增强规划器的可行性恢复保证，证明最大间隙空间决策的$O(n\log n)$最优性；
  \item 在Apollo 11.0生产级C++代码上实现框架并完成实车验证，以最小侵入度（$\sim$800行新增代码）实现压力场景通过率从0\%到100\%的提升。
\end{enumerate}

% ============================================================
\section{问题建模与解耦失效分析}
% ============================================================

\subsection{解耦规划架构形式化}

考虑自车沿参考线以Frenet坐标$(s, l)$描述。解耦规划将轨迹生成分解为两个顺序子问题：

\textbf{路径子问题（Stage 1）：} 给定道路边界$[l_{\min}(s), l_{\max}(s)]$和障碍物占用区域$\mathcal{B}(s)$，求解横向偏移曲线$l^*(s)$使得
\begin{equation}
  l^* = \arg\min_{l} J_{\text{path}}(l) \quad \text{s.t.} \quad l(s) \in [l_{\min}(s), l_{\max}(s)] \setminus \mathcal{B}(s), \;\forall s
\end{equation}
其中$J_{\text{path}}$包含平滑性、参考线跟踪等代价项。关键在于$\mathcal{B}(s)$的构造：在纯解耦架构中，动态障碍物$o_i$的占用通过对所有预测时刻取并集得到：
\begin{equation}\label{eq:static-proj}
  \mathcal{B}_{\text{dec}}(s) = \bigcup_{t=0}^{T_{\text{pred}}} \text{Proj}_{SL}(o_i, t)
\end{equation}
这种"全时刻并集"投影丢失了时间信息，将动态障碍物等效为其整个预测轨迹的空间包络。

\textbf{速度子问题（Stage 2）：} 给定固定路径$l^*(s)$和ST域边界约束，求解速度曲线$s^*(t)$使得
\begin{equation}
  s^* = \arg\min_{s} J_{\text{speed}}(s) \quad \text{s.t.} \quad (s(t), t) \in \mathcal{S}_{\text{free}}, \;\forall t
\end{equation}
其中$\mathcal{S}_{\text{free}}$为ST图中的自由空间。

两阶段的唯一接口为路径$l^*(s)$——一条不含时间信息的空间曲线。

\subsection{信息损失与失效机制}

定义联合可行域$\mathcal{F}_{\text{joint}}$为同时满足所有空间和时序避障约束的轨迹集合，解耦可行域$\mathcal{F}_{\text{dec}}$为经式(\ref{eq:static-proj})投影后两阶段顺序求解可达的轨迹集合。由于静态投影是时序占用的超集，有：
\begin{equation}
  \mathcal{F}_{\text{dec}} \subseteq \mathcal{F}_{\text{joint}}
\end{equation}

当动态障碍物的预测包络足够大时，$\mathcal{F}_{\text{dec}}$可退化为空集而$\mathcal{F}_{\text{joint}}$仍非空。这就是解耦规划的结构性失效。

\subsection{解耦不可行性条件}

\textbf{命题1（解耦不可行性）：} 设动态障碍物$o_i$的预测轨迹为$\hat{p}_i(t)$，自车道路横向可行域为$[l_{\min}(s), l_{\max}(s)]$。若存在纵向区间$[s_a, s_b]$使得
\begin{equation}
  \forall s \in [s_a, s_b]: \quad \bigcup_{t=0}^{T_{\text{pred}}} \text{Proj}_{SL}(o_i, t) \supseteq [l_{\min}(s) + W/2,\; l_{\max}(s) - W/2]
\end{equation}
其中$W$为自车宽度，则纯解耦规划器的路径子问题在$[s_a, s_b]$上无可行解，即$\mathcal{F}_{\text{dec}} = \emptyset$。

\textit{证明.} 路径子问题依据式(\ref{eq:static-proj})构造占用区域，不具备查询特定时刻占用的能力。因此其观测到的可行带宽度为
$g_{\text{dec}}(s) = (l_{\max}(s) - l_{\min}(s)) - |\bigcup_t \text{Proj}_{SL}(o_i, t)|$。
当上述条件成立时，$g_{\text{dec}}(s) < W$，$\forall s \in [s_a, s_b]$，
即无法为自车分配宽度为$W$的无碰撞通行带。路径子问题报告 BLOCKED，规划终止。

然而，若规划器知道自车到达$s$的时刻$\tau(s)$，则只需查询$o_i$在$t = \tau(s)$时刻的占用：
\begin{equation}
  \mathcal{B}_{\text{enrich}}(s) = \text{Proj}_{SL}(o_i, \tau(s))
\end{equation}
这通常远小于全时刻并集，可行带宽度$g_{\text{enrich}}(s) > W$的概率显著增加。\hfill$\square$

这一命题揭示了解耦失效的根本原因：时间维度信息的缺失导致占用区域被过度膨胀。由此引出本文的核心思路——将时间信息以约束形式传递给速度阶段，等效于在路径阶段获得时间感知能力。

% ============================================================
\section{约束增强解耦规划框架}
% ============================================================

本节按信息流方向阐述CEDP框架的设计。框架在解耦架构的两阶段接口处插入一条约束传递通道$\mathcal{T}$，其携带结构化信息$\mathcal{T} = \{(s_k, \tau_k, \delta_k)\}_{k=1}^{K}$，分别表示关键纵向位置、自车预计到达时间和差异化安全裕度。通道的构造由三个协同组件完成。

\subsection{框架总体架构}

CEDP的信息流如下：
\begin{equation}
\underbrace{\text{感知} \to \text{威胁度过滤} \to \text{空间决策} \to \text{路径QP}}_{\text{Stage 1（增强）}}
\xrightarrow{\;\mathcal{T}\;}
\underbrace{\text{走廊构造} \to \text{速度QP}}_{\text{Stage 2（增强）}}
\end{equation}

与原始解耦架构相比，唯一的结构性变化是通道$\mathcal{T}$的引入。Stage 1和Stage 2内部的QP求解器保持不变，仅约束集合被通道信息扩充。

\subsection{威胁度过滤：确定传递对象}

并非所有障碍物都需要进入约束通道。静态障碍物和远离自车的低威胁目标可由原始解耦架构正常处理。通道仅服务于可能触发解耦失效的高威胁动态障碍物。

定义障碍物$o_i$的综合威胁度：
\begin{equation}
  \Theta_i = \sum_{j=1}^{5} w_j f_j(o_i)
\end{equation}
其中五个因子分别为：
\begin{itemize}
  \item $f_{\text{TTC}}$：碰撞时间因子，$f_{\text{TTC}} = \max(0, 1 - \text{TTC}_i / T_{\text{crit}})$，$T_{\text{crit}} = 1.5$\,s；
  \item $f_{\text{overlap}}$：横向重叠比，自车与障碍物在SL域的几何重叠程度；
  \item $f_{\text{vel}}$：相对速度因子，Sigmoid映射的接近速度；
  \item $f_{\text{type}}$：障碍物类型因子，行人/自行车$= 0.8$，车辆$= 0.5$，静态$= 0.1$；
  \item $f_{\text{density}}$：交互密度因子，邻近障碍物的聚集程度。
\end{itemize}

仅$\Theta_i > \Theta_{\text{thr}}$的障碍物进入后续空间决策环节。该过滤器将通道负载从全部$N$个障碍物降至$K \ll N$个高威胁目标。

\subsection{空间决策：最大间隙通行带选择}

对通过威胁度过滤的$K$个障碍物，需要确定自车的横向通行策略。在多障碍物场景下，Apollo原始的逐障碍物贪心nudge决策可能产生矛盾方向（左绕$o_1$与右绕$o_2$冲突），导致路径子问题无解。

本文采用扫描线分组与最大间隙策略：

\textbf{步骤1：纵向分组。} 对$K$个障碍物按纵向占用区间$[s_i^{\text{start}}, s_i^{\text{end}}]$进行扫描线分组，将纵向重叠的障碍物归入同一组$G_m$。

\textbf{步骤2：组内间隙计算。} 对每组$G_m$中的$k$个障碍物，按横向中心$u_i$排序，计算$k+1$个候选间隙的宽度。

\textbf{步骤3：最大间隙选取。}

\textbf{定理（最大间隙最优性）：} 设组$G_m$内有$k$个障碍物，其横向边界排序后为$\{[u_i, v_i]\}_{i=1}^{k}$（$u_i \leq u_{i+1}$）。定义$V_p = \max_{q \leq p} v_q$为前$p$个障碍物右边界的运行最大值。则最优通行带宽度为
\begin{equation}
  g^* = \max_{0 \leq p \leq k} g_p
\end{equation}
其中
\begin{equation}
  g_p = \begin{cases}
    u_1 - l_{\min} & p = 0 \\
    u_{p+1} - V_p & 1 \leq p \leq k-1 \\
    l_{\max} - V_k & p = k
  \end{cases}
\end{equation}

\textit{证明.} 首先证明引理（连续划分性）：最优通行带必然对应一个连续的左/右划分，即存在分割点$p^*$使得障碍物$1, \ldots, p^*$位于通行带右侧，$p^*+1, \ldots, k$位于左侧。这由交换论证得到——若最优解中存在非连续划分，可通过交换相邻障碍物的绕行方向获得不劣的连续划分。

基于引理，最优通行带宽度等于$k+1$个候选间隙中的最大值。排序需$O(k \log k)$，扫描需$O(k)$，总复杂度$O(k \log k)$。对所有组求和，总复杂度为$O(n \log n)$。\hfill$\square$

空间决策输出最优通行带边界$[l_{\min}^*(s), l_{\max}^*(s)]$，送入路径QP作为约束。

\subsection{时间约束注入：时空可通行走廊}

路径QP求解得到$l^*(s)$后，结合自车运动学模型计算到达时间函数：
\begin{equation}
  \tau(s) = \int_0^s \frac{1}{v(\sigma)} d\sigma
\end{equation}
其中$v(\sigma)$为基于当前速度和加速度能力的预估速度剖面。

对每个高威胁障碍物$o_i$，查询其在$t = \tau(s)$时刻的SL投影，构造时变通行带：
\begin{equation}
  g_{\text{TV}}(s) = (l_{\max}(s) - l_{\min}(s)) - |\text{Proj}_{SL}(o_i, \tau(s))|
\end{equation}

定义时空可通行走廊：
\begin{equation}
  \Omega = \{(s, t) \mid g_{\text{TV}}(s) \geq W,\; t \geq \tau(s)\}
\end{equation}

将$\Omega$转化为速度QP的线性约束后注入Stage 2，与原有SpeedBoundsDecider的约束取交集。

\textbf{命题2（可行性恢复）：} 若$\Omega \neq \emptyset$且在速度QP的离散化网格上构成凸集，则增强后的速度QP必有可行解。

\textit{证明.} 速度QP的目标函数为连续二次函数（凸），约束集为$\Omega$与原有ST边界的交集。当$\Omega \neq \emptyset$且凸时，交集非空且凸，凸QP在非空凸可行域上必有最优解。\hfill$\square$

\subsection{验证反馈环}

速度QP求解得到$s^*(t)$后，执行一次一致性检查：验证实际速度曲线隐含的到达时间$\tau'(s) = (s^*)^{-1}(s)$是否与构造走廊时使用的$\tau(s)$一致。若偏差超过阈值$\epsilon_\tau$：
\begin{equation}
  |\tau'(s) - \tau(s)| > \epsilon_\tau, \quad \exists s \in [s_a, s_b]
\end{equation}
则用$\tau'(s)$重新裁剪走廊$\Omega'$并再次求解速度QP。该反馈最多执行一次，计算代价为额外一次QP求解（$< 5$\,ms）。

% ============================================================
\section{Apollo 11.0 实现}
% ============================================================

CEDP在Apollo 11.0的PublicRoad Planner流水线中实现，修改涉及三个模块：

\begin{itemize}
  \item \texttt{PathBoundsDecider}：注入威胁度计算和最大间隙决策逻辑，替换原始逐障碍物贪心nudge；
  \item \texttt{SpeedBoundsDecider}：接收走廊约束$\Omega$，与原有ST边界取交集后送入速度QP；
  \item 新增\texttt{ConstraintChannel}：管理通道数据结构$\mathcal{T}$的生命周期。
\end{itemize}

实现统计：新增C++代码约800行，未修改任何QP求解器内部逻辑（仍使用OSQP），未引入新的外部依赖。框架与Apollo原有的fallback机制（LaneBorrowPath、EmergencyStop）保持兼容。

% ============================================================
\section{实验验证}
% ============================================================

\subsection{实验设置}

\textbf{仿真平台：} Apollo 11.0 Dreamview，规划频率10\,Hz，车辆运动学模型。

\textbf{场景设计：} 15个场景分为三类：
\begin{itemize}
  \item 压力场景（P1--P6）：多动态障碍物交互，纯解耦架构必然失效；
  \item 可比场景（C1--C6）：中等复杂度，解耦架构可通过但性能受限；
  \item 边界场景（B1--B3）：参数极端化，测试CEDP的失败边界。
\end{itemize}

每个场景设置5组参数变体（障碍物速度$\pm 20\%$、$\pm 40\%$、障碍物数量$\pm 1$），共计$15 \times 5 = 75$组配置。

\textbf{对比方法：}
\begin{itemize}
  \item Baseline：原始Apollo EM Planner（纯解耦）；
  \item CEDP：本文方法；
  \item CEDP-ablation：分别关闭威胁度过滤/间隙决策/走廊注入/反馈环。
\end{itemize}

\textbf{评价指标：} 通过率、平均巡航速度、最小障碍物距离、纵向jerk、规划单帧耗时。

\subsection{仿真结果}

（此处待补充实验数据表格）

\subsection{消融实验}

（此处待补充消融实验结果）

\subsection{计算开销分析}

（此处待补充全流水线延迟分解表）

\subsection{失败边界分析}

（此处待补充参数敏感性曲线：在什么障碍物密度/速度下CEDP也失效）

\subsection{实车验证}

\textbf{平台：} （待补充车辆型号、传感器配置、计算硬件）

\textbf{场景：} 封闭场地，2--3个锥桶/假人布置，覆盖静态绕行和动态交互。

\textbf{结果：} （待补充规划指令曲线、实际执行曲线、场景截图）

% ============================================================
\section{相关工作}
% ============================================================

\subsection{时空联合规划}

IA Planner通过瞬时分析将时空约束降维后用MPC联合求解。胡杰等人在$s$-$l$-$t$空间中用Hybrid A*搜索初始轨迹再用NMPC优化。乔应基于三维DP搜索配合AABB走廊约束。这类方法彻底放弃解耦架构，以计算代价翻倍为代价获得时空一致性。本文的CEDP保留解耦架构，通过通道传递实现近似等效的时空协调，计算代价仅增加一次QP。

\subsection{解耦框架内改进}

Han等人在路径QP目标函数中注入多模态预测概率作为soft cost，隐式传递预测信息，但缺乏可行性保证。Zhang等人通过并行枚举路径--速度组合并用合作博弈选取最优配对，复杂度为$O(mn)$。Tong等人提出轨迹修复框架，在解耦失效时升级为时空联合修复。本文与上述工作的区别在于：以硬约束形式显式传递结构化信息，具有可行性保证（命题2），且不改变各阶段内部求解结构。

\subsection{解耦规划的工业应用}

Meng等人提出Lattice+凸优化的解耦框架。Wang等人报告了京东配送车的工业级解耦规划系统。这些工作验证了解耦架构在工业部署中的主导地位，也间接说明了在该架构内做改进的实际价值。

% ============================================================
\section{结论}
% ============================================================

本文提出约束增强解耦规划框架CEDP，通过在路径阶段与速度阶段之间建立显式约束传递通道，在保留解耦架构计算优势的前提下恢复时空协调能力。理论分析表明纯解耦架构在动态障碍物覆盖条件下必然失效，而CEDP在走廊非空时保证可行性。在Apollo 11.0生产级代码上的实现和实车验证表明，框架以最小侵入度实现了压力场景通过率的显著提升。

未来工作将在两个方向展开：一是将约束通道扩展为双向反馈（速度信息回传路径阶段），实现更紧密的迭代耦合；二是在大规模随机场景下进行统计验证，量化CEDP相对于联合规划的次优性上界。

\end{document}
